
\phantomsection\addcontentsline{toc}{chapter}{Laboratorio matemático: de conteos a demostraciones}\label{lab:mat}
\chapter*{Laboratorio matemático de pizarra: estimar y comprobar}
\addcontentsline{toc}{chapter}{Laboratorio matemático de pizarra}

\section*{Un experimento completo desde cero: cuatro símbolos en prueba}
\addcontentsline{toc}{section}{Un experimento completo desde cero}
Fijemos el alfabeto $V=\{a,b,c\}$, una cadena de entrenamiento
$\texttt{abacabbacb}$ y la cadena reservada $\texttt{abca}$.
Ambas son \textbf{ILUSTRATIVAS}: se construyeron para enseñar la cuenta, no provienen de
un corpus observado. No obstante, sus resultados numéricos están regenerados por
\texttt{codigo/s01.py} y contrastados en \texttt{verificaciones/}.

Construimos la matriz de transiciones $C_{ij}=\#\{t:x_t=i,x_{t+1}=j\}$:
\[
C=\begin{array}{c|rrr|r}&a&b&c&N_i\\\hline
 a&0&2&2&4\\ b&2&1&0&3\\ c&1&1&0&2\end{array}.
\]
Hay $10-1=9$ transiciones, no diez: el símbolo inicial no tiene predecesor.
La ecuación de normalización para $\alpha=1$ es
\[
 P_{ij}=\frac{C_{ij}+1}{N_i+3}
 =\begin{pmatrix}1/7&3/7&3/7\\1/2&1/3&1/6\\2/5&2/5&1/5\end{pmatrix}.
\]
El denominador de la primera fila es $7$, no $6$; el de la tercera es $5$.
La comprobación de suma por fila es elemental pero obligatoria:
$1/7+3/7+3/7=1$, $1/2+1/3+1/6=1$ y $2/5+2/5+1/5=1$.
La probabilidad \emph{condicional de las tres transiciones de prueba}, dado $X_1=a$, es
\[
 P(b\mid a)P(c\mid b)P(a\mid c)
 =\frac37\cdot\frac16\cdot\frac25=\frac1{35}.
\]
No es todavía la probabilidad de la cadena completa. Para ello hay que precisar
la distribución inicial $\pi$: si elegimos $\pi=(1/3,1/3,1/3)$,
\[
 q(\texttt{abca})=\pi(a)\,P(b\mid a)P(c\mid b)P(a\mid c)=\frac1{105}.
\]
Por tanto, la longitud ideal de código es $\log_2(105)\approx6.714246$ bits,
y, con cuatro eventos de predicción, la tasa es
$\log_2(105)/4\approx1.678561$ bits por carácter y la perplejidad
$105^{1/4}\approx3.201086$.
Si en lugar de cuatro eventos se evalúan solo las tres transiciones,
la pérdida media pasa a $\log_2(35)/3$; \emph{ambas cifras son válidas,
pero corresponden a experimentos condicionales diferentes}.

\begin{ejemplo}[Un contraste que detecta un modelo incorrecto]
El estimador no suavizado asigna $\widehat P(c\mid b)=0$.
Entonces $q(\texttt{abca})=0$ para cualquier $\pi(a)>0$ y la pérdida
es $+\infty$. Sustituir el cero únicamente al mirar la prueba sería
modificar el modelo después de observarla: hay que definir el suavizado y
seleccionar su intensidad en validación.
\end{ejemplo}

\begin{figure}[H]\centering
\begin{tikzpicture}[>=Stealth,font=\small, every node/.style={align=center}]
\node[draw=azul,fill=azulclaro,rounded corners,minimum width=3cm,minimum height=1cm] (t) at (0,0) {\textbf{Entrenamiento}\\10 símbolos};
\node[draw=azul,fill=azulclaro,rounded corners,minimum width=3cm,minimum height=1cm] (c) at (4,0) {\textbf{Conteos}\\9 transiciones};
\node[draw=verde,fill=verdeclaro,rounded corners,minimum width=3cm,minimum height=1cm] (p) at (8,0) {\textbf{Suavizado}\\filas de $P$};
\node[draw=ocre,fill=ocreclaro,rounded corners,minimum width=3cm,minimum height=1cm] (e) at (8,-2) {\textbf{Prueba}\\$\texttt{abca}$};
\node[draw=azul,fill=azulclaro,rounded corners,minimum width=3cm,minimum height=1cm] (r) at (2,-2) {\textbf{Evaluación}\\$\log_2 105$ bits};
\draw[->,thick] (t)--(c);\draw[->,thick] (c)--(p);\draw[->,thick] (p)--(e);\draw[->,thick] (e)--(r);
\end{tikzpicture}
\caption{La prueba solo recibe un modelo previamente fijado. El flujo contrario constituye fuga de información.}
\end{figure}

\section*{Regla de la cadena frente a memoria finita: una cuenta contrastable}
\addcontentsline{toc}{section}{Regla de la cadena y memoria finita}
Tomemos $X_1$ y $X_2$ Bernoulli$(1/2)$ independientes, y definamos
$X_3=X_1$ casi seguramente. Se trata de una distribución conjunta legítima,
aunque no sea markoviana de orden uno. Para $x=(1,0,1)$,
\[
 p(1,0,1)=p(X_1=1)\,p(X_2=0\mid X_1=1)\,
 p(X_3=1\mid X_1=1,X_2=0)=\frac12\frac12\cdot1=\frac14.
\]
Un modelo construido con los \emph{verdaderos} condicionales de orden uno
utilizaría $p(X_3=1\mid X_2=0)=1/2$ y asignaría $1/8$ a esa cadena.
Esta diferencia no es error algebraico: proviene de descartar información.

\begin{figure}[H]\centering
\begin{tikzpicture}[>=Stealth,font=\small,
 v/.style={draw=azul,fill=azulclaro,circle,minimum size=10mm}]
\node[v] (x1) at (0,0) {$X_1$};\node[v] (x2) at (4,0) {$X_2$};\node[v] (x3) at (8,0) {$X_3$};
\draw[->,thick] (x1)--(x2);\draw[->,thick] (x2)--(x3);
\draw[->,thick,dashed,color=vino] (x1) to[bend left=35] node[above] {memoria relevante} (x3);
\node at (4,-1.0) {Un bigrama descarta la dependencia discontinua $X_1\rightarrow X_3$.};
\end{tikzpicture}
\caption{Diagrama explicativo de una dependencia que no cabe en un único símbolo de memoria.}
\end{figure}

\section*{Suavizado como regularización: el papel del tamaño del vocabulario}
\addcontentsline{toc}{section}{Suavizado y tamaño del vocabulario}
Sea $K$ el tamaño del alfabeto y $\alpha>0$. El peso total de los
pseudoconteos es $K\alpha$, de modo que
\[
\frac{c_i+\alpha}{N+K\alpha}
=\underbrace{\frac{N}{N+K\alpha}}_{\lambda}\frac{c_i}{N}
 +\underbrace{\frac{K\alpha}{N+K\alpha}}_{1-\lambda}\frac1K.
\]
Esta identidad describe el suavizado como mezcla convexa entre el estimador
empírico y la distribución uniforme. Es válida para $N>0$ y $\alpha>0$.
Para $N=10$, $K=5$ y $\alpha=1$, $\lambda=2/3$; si se aumenta
$K$ a 100, manteniendo $N$ y $\alpha$, $\lambda=1/11$. La masa que
se introduce por suavizado depende fuertemente del vocabulario.

\begin{ejemplo}[Una decisión que exige validación]
En un contexto con conteos $(6,3,1,0,0)$, $K=5$ y $\alpha=1$,
se obtiene $q=(7,4,2,1,1)/15$. El símbolo ausente recibe $1/15$.
Con $\alpha=1/2$, el denominador es $12.5$ y recibe $0.5/12.5=1/25$.
Ninguna de las dos opciones es automáticamente superior: la pérdida en
validación sirve para decidir antes de examinar prueba.
\end{ejemplo}

\section*{Perplejidad: tres lecturas compatibles y una lectura incorrecta}
\addcontentsline{toc}{section}{Perplejidad y medias geométricas}
Fijadas las probabilidades $r_t=q(x_t\mid h_t)>0$ de $T$ eventos observados,
\[
\operatorname{PPL}
=\exp\left(-\frac1T\sum_{t=1}^T\ln r_t\right)
=\left(\prod_{t=1}^T\frac1{r_t}\right)^{1/T}.
\]
En consecuencia, la PPL es la \emph{media geométrica de las inversas}
de las probabilidades asignadas a los eventos. Si todas son iguales a
$1/K$, la PPL es $K$; para probabilidades variables, la PPL no es
el número real de símbolos candidatos ni la inversa de la exactitud.

\begin{ejemplo}[Producto, bits y exponencial: cálculo completo]
Para probabilidades observadas $(1/2,1/4,1/8,1/2)$,
\[
-\sum_{t=1}^4\log_2r_t=1+2+3+1=7\text{ bits},
\qquad\operatorname{PPL}=2^{7/4}\approx3.363586.
\]
Si la tercera probabilidad se reduce de $1/8$ a $1/16$ sin alterar las demás,
la suma aumenta a $8$ bits y la PPL a $2^2=4$. El cambio puede
explicarse multiplicativamente: se duplicó el código ideal de ese evento.
\end{ejemplo}

\section*{Comparar tasas, no solo nombres de métricas}
\addcontentsline{toc}{section}{Comparación justa entre tokenizaciones}
Dos modelos evalúan exactamente la misma cadena de $100$ caracteres,
sin pérdida de información al tokenizar. El modelo $A$ predice
$100$ caracteres con PPL $3.2$, mientras que $B$ predice $30$ subpalabras
con PPL $25$. Sus longitudes ideales son
\[
B_A=100\log_2(3.2)\approx167.8072\text{ bits},\quad
B_B=30\log_2(25)\approx139.3157\text{ bits}.
\]
Si el texto de referencia y el cálculo probabilístico sobre \emph{todo} el texto
son comparables, las tasas respectivas son $1.6781$ y $1.3932$ bits por
carácter. En este ejemplo el orden por PPL es inverso al orden por tasa.

\begin{nota}[Limitación esencial]
La misma unidad de longitud no convierte esta comparación en un efecto
\emph{causal} de la tokenización: también pueden variar datos de entrenamiento,
normalización, parametrización, contexto y presupuesto computacional. En particular,
no se identifica con BPC la calidad semántica de un modelo.
\end{nota}
Tampoco debe contarse el mismo dato dos veces: si una cadena Unicode se
normaliza antes del modelado, el número de caracteres del denominador debe
reflejar una convención fija y documentada. Para comparar código sin pérdida
se requiere poder reconstruir exactamente la cadena original.

\section*{De muestras a poblaciones: qué se demuestra realmente}
\addcontentsline{toc}{section}{Diferencia entre riesgo empírico y poblacional}
La pérdida observada sobre $T$ eventos es
$\widehat R_T(q)=T^{-1}\sum_{t=1}^T-\ln q(x_t\mid h_t)$.
En un problema categórico \emph{i.i.d.} con observaciones de distribución
$p$ y soporte compatible, la esperanza de $-\ln q(X)$ es $H(p,q)$;
la ley de los grandes números justifica convergencia para $q$ fijo
si la variable es integrable. En una secuencia dependiente la convergencia
requiere hipótesis apropiadas (por ejemplo, ergodicidad e integrabilidad),
y en modelos elegidos usando los mismos datos aparece un problema adicional
de sobreajuste.

Por otro lado, para una muestra fija, minimizar $\widehat R_T(q_\theta)$
es exactamente maximizar la verosimilitud correspondiente, porque
\[
\arg\min_\theta\widehat R_T(q_\theta)
=\arg\min_\theta\left(-\frac1T\ln q_\theta(x_{1:T})\right)
=\arg\max_\theta q_\theta(x_{1:T}).
\]
Esto no prueba que $q_\theta=p$, ni establece por sí solo generalización.
El paso de lo observado a una afirmación científica exige particiones,
incertidumbre y una pregunta experimental bien definida.


\chapter*{Derivaciones complementarias y experimentos de pizarra}
\addcontentsline{toc}{chapter}{Derivaciones complementarias y experimentos de pizarra}

\section*{De una frase a una variable aleatoria}
\addcontentsline{toc}{section}{De una frase a una variable aleatoria}
Antes de introducir redes neuronales separamos tres operaciones distintas:
(i) elegir una representación discreta del texto, (ii) asignar probabilidades a
secuencias en esa representación y (iii) generar símbolos condicionando en un
prefijo. Por ejemplo, la cadena ``pozo estable'' puede representarse por caracteres,
palabras o subpalabras; las probabilidades sobre los elementos de cada representación
no tienen por qué coincidir. Una representación admite un modelo autorregresivo
si permite definir $q(x_t\mid x_{<t})$ y su producto a lo largo de la secuencia.

\begin{figure}[H]\centering
\begin{tikzpicture}[font=\small,box/.style={draw=azul,rounded corners=2pt,fill=azulclaro,
minimum height=9mm,align=center,text width=3cm},>=Stealth]
\node[box] (a) at (0,0) {Texto original\\\texttt{pozo estable}};
\node[box] (b) at (4.1,0) {Tokenización\\$x_1,\ldots,x_T$};
\node[box] (c) at (8.2,0) {Probabilidad\\$\prod_tq(x_t\mid x_{<t})$};
\draw[->,thick] (a)--(b);
\draw[->,thick] (b)--(c);
\node[draw=gris,dashed,rounded corners,text width=7.7cm,align=center] at (4.1,-1.5)
{Medición: $-\log_2 q(x_{1:T})$ bits; fijar fronteras y unidad del texto};
\end{tikzpicture}
\caption{Tres decisiones separadas: representación, asignación de probabilidad y medición.}
\end{figure}

\section*{Regla de la cadena con números: dependencia que el bigrama pierde}
\addcontentsline{toc}{section}{Regla de la cadena con números}
Sea $X_1\sim\mathrm{Bernoulli}(1/2)$, $X_2\sim\mathrm{Bernoulli}(1/2)$ independiente
y $X_3=X_1$ con probabilidad uno. La regla de la cadena es exacta:
$p(x_1,x_2,x_3)=p(x_1)p(x_2\mid x_1)p(x_3\mid x_1,x_2)$.
La probabilidad de $(1,0,1)$ es $1/2\cdot1/2\cdot1=1/4$. Sin embargo,
$p(X_3=1\mid X_2=0)=1/2$; al imponer el supuesto de bigramas se obtiene
$1/2\cdot1/2\cdot1/2=1/8$. La diferencia se debe a la memoria perdida, no
a una falla de la regla de la cadena. En esta distribución el pasado remoto
$X_1$ es indispensable para predecir $X_3$.

\section*{La estimación por máxima verosimilitud cuando hay ceros}
\addcontentsline{toc}{section}{MLE y sus casos de frontera}
Para un contexto $h$, sean $c_1,\ldots,c_K$ los conteos y $N=\sum_i c_i>0$.
El objetivo $\ell(q)=\sum_{i:c_i>0}c_i\ln q_i$ se maximiza sobre el símplex
cerrado $q_i\geq0$, $\sum_iq_i=1$. Sobre el soporte observado, las condiciones
de Lagrange dan $c_i/q_i=\lambda$. Por ello $q_i=c_i/N$, y por el principio
de la desigualdad de Gibbs este valor maximiza globalmente la log-verosimilitud.
Para un $c_i=0$, la solución asigna $q_i=0$; una solución estrictamente positiva
para todas las categorías no alcanza el máximo, solo puede aproximarlo.

\begin{center}
\begin{tabular}{lrrrr}\toprule
Contexto $h$ & $c(a)$ & $c(b)$ & $c(c)$ & $N_h$\\\midrule
$a$ & 0 & 2 & 2 & 4\\
$b$ & 2 & 1 & 0 & 3\\
$c$ & 1 & 1 & 0 & 2\\\bottomrule
\end{tabular}
\end{center}
Para el contexto $b$, $\hat q(c\mid b)=0$. Con $\alpha=1/2$ y $K=3$,
\[
q_\alpha(\cdot\mid b)=\biggl(\frac{2.5}{4.5},\frac{1.5}{4.5},
\frac{0.5}{4.5}\biggr)=\left(\frac59,\frac13,\frac19\right).
\]
La masa para el símbolo no visto es $1/9$; para $\alpha=1$ sería $1/6$.
Esto pone de relieve que el parámetro de suavizado cambia predicciones,
por lo que se selecciona en validación, nunca en prueba.

\section*{De Dirichlet a predicción posterior}
\addcontentsline{toc}{section}{De Dirichlet a predicción posterior}
Sea $\mathbf q\sim\operatorname{Dirichlet}(\alpha_1,\ldots,\alpha_K)$.
Su densidad (hasta constante) es $\prod_iq_i^{\alpha_i-1}$.
Dada una muestra con conteos $\mathbf c$, la verosimilitud es proporcional
a $\prod_iq_i^{c_i}$, de modo que la posterior tiene densidad
$\prod_iq_i^{c_i+\alpha_i-1}$; por identificación,
$\mathbf q\mid \mathbf c\sim\operatorname{Dirichlet}(\mathbf c+\boldsymbol\alpha)$.
La distribución predictiva para una nueva observación es
\[
P(X_{N+1}=i\mid\mathbf c)=\int q_i\,p(\mathbf q\mid\mathbf c)\,d\mathbf q
=\frac{c_i+\alpha_i}{N+\sum_j\alpha_j}.
\]
La integral es una esperanza posterior, no una moda posterior ni un estimador
frecuentista sin regularización. Este punto facilita comprender que un modelo
puede conservar incertidumbre sobre parámetros a partir de pocos datos.

\section*{Entropía, cruzada, divergencia y codificación}
\addcontentsline{toc}{section}{Entropía, cruzada, divergencia y codificación}
Para $p=(3/4,1/4)$ y $q=(1/2,1/2)$, usando logaritmos en base dos,
\[
H(p)=-\tfrac34\log_2\tfrac34-\tfrac14\log_2\tfrac14\simeq0.8113,
\qquad H(p,q)=1.
\]
Luego $D_{\mathrm{KL}}(p\Vert q)=H(p,q)-H(p)\simeq0.1887$ bits.
Invirtiendo el orden se obtiene
$D_{\mathrm{KL}}(q\Vert p)=\frac12\log_2(2/3)+\frac12\log_2 2
\simeq0.2075$ bits: KL no es simétrica.
El argumento de Gibbs establece $D_{\mathrm{KL}}(p\Vert q)\geq0$;
si $q=0$ en un punto con $p>0$, se entiende que la divergencia es infinita.
La interpretación de codificación es operacional: $-\log_2 q(x)$ mide
longitud idealizada de un código probabilístico para el evento $x$.

\section*{Tres modelos, una cadena: por qué la unidad importa}
\addcontentsline{toc}{section}{Tres modelos, una cadena}
Considérese un texto de $100$ caracteres Unicode medidos como puntos de
código, sin normalización posterior. Supónganse dos distribuciones válidas
sobre sus secuencias de tokens, con codificación reversible del texto:
\begin{center}
\begin{tabular}{lrrrr}\toprule
Modelo & Tokens & PPL/token & Bits totales & BPC\\\midrule
A (caracteres) &100& $3.2$ & $100\log_2(3.2)\simeq167.8$ & $1.678$\\
B (subpalabras)&30& $25$ & $30\log_2(25)\simeq139.3$ & $1.393$\\\bottomrule
\end{tabular}
\end{center}
Ambos juegos de probabilidades son \emph{ilustrativos}. Para comparar tasas
sobre un mismo texto debe incluirse toda la probabilidad del evento textual,
incluidos marcadores de inicio y fin cuando correspondan, y garantizar que
ambas representaciones reconstruyan el texto sin pérdida. Si un tokenizador
normaliza y destruye caracteres, el BPC calculado sobre texto original
puede dejar de representar una tasa comparable; es necesario especificar
la transformación o evaluarla explícitamente.

\section*{Ventana de contexto y mecanismo de generación}
\addcontentsline{toc}{section}{Ventana de contexto y mecanismo de generación}
En el generador autorregresivo, el \emph{prompt} proporciona el prefijo;
el contexto es la historia efectivamente suministrada al modelo;
y el arnés decide si muestrear, terminar o truncar la respuesta.
\begin{figure}[H]\centering
\begin{tikzpicture}[font=\small,>=Stealth,
box/.style={draw=azul,rounded corners,fill=azulclaro,align=center,text width=3cm,minimum height=1cm}]
\node[box] (p) at (0,0) {Prefijo\,(prompt)\\$x_{1:k}$};
\node[box] (q) at (4,0) {Distribución condicional\\$q_\theta(\cdot\mid h)$};
\node[box] (t) at (8,0) {Muestreo de $X_{k+1}$\\+ condición de parada};
\draw[->,thick] (p)--(q);
\draw[->,thick] (q)--(t);
\draw[->,thick] (t.south) -- ++(0,-.7) -| node[pos=.25,below]{si continúa: actualizar $h$} (p.south);
\end{tikzpicture}
\caption{El modelo calcula condicionales; el arnés transforma condicionales en un algoritmo de generación.}
\end{figure}
En particular, una cota de $n$ tokens impuesta externamente produce
una distribución de longitudes diferente de la ley generativa sin truncamiento.

\section*{Puente hacia la práctica en R y Python}
\addcontentsline{toc}{section}{Puente hacia la práctica en R y Python}
La estrategia pedagógica separa la representación de texto, la construcción
de una matriz de conteos y su interpretación probabilística. El módulo
\texttt{modulo\_codigo\_s01.tex} proporciona un corpus ilustrativo, código base R
sin dependencias, validación de sumas por fila, cálculo de entropía cruzada
y una implementación didáctica de BPE con desempates deterministas.
Las arquitecturas recurrentes, convolucionales y \emph{Transformer} se
mencionan solo como extensiones: su tratamiento detallado corresponde
a sesiones posteriores.

\chapter*{Cierre del laboratorio matemático: predicciones que el código debe verificar}
\phantomsection\addcontentsline{toc}{chapter}{Cierre del laboratorio matemático}\label{lab:mat:cierre}
\begin{enumerate}
\item Sobre \texttt{abacabbacb}, la fila suavizada de $a$ es
$\left(\frac17,\frac37,\frac37\right)$, no $(1/6,\ldots)$:
$N_a=4$, $K=3$ y $\alpha=1$.
\item Para \texttt{abca}, la \emph{probabilidad condicionada en el primer
símbolo} es el producto de tres transiciones. La probabilidad \emph{conjunta}
incluye la masa inicial declarada. Mezclar ambos protocolos cambia la PPL.
\item La identidad $H(p,q)=H(p)+D_{\mathrm{KL}}(p\Vert q)$ es una igualdad
matemática (con convenciones de soporte), no una afirmación de que un modelo
entrenado estima correctamente $p$.
\item En una evaluación reproducible, la unidad de resampling es el documento;
no se deben remuestrear caracteres relacionados como observaciones independientes.
\end{enumerate}
\begin{puente}[Puente al laboratorio computacional]
La teoría define qué se calcula. En el laboratorio computacional de esta sesión
se declara el corpus, se fijan reglas de preprocesamiento y se generan
\texttt{resultados\_python.csv} y \texttt{resultados\_R.csv}; la concordancia
numérica de ambos es una prueba de implementación, no de generalización.
\end{puente}
